\documentclass[UTF8]{ctexart}
\usepackage{xeCJK}
\usepackage{geometry}
\geometry{left=2cm,right=2cm,top=2.5cm,bottom=2.5cm}
\setlength{\headheight}{12.7pt}

\usepackage{titlesec}
\titleformat{\section}
  {\normalfont\large\bfseries}
  {\thesection}
  {1em}
  {}
\titlespacing*{\section}
  {0pt}
  {3.5ex plus 1ex minus .2ex}
  {2.3ex plus .2ex}

\usepackage{amsmath}
\usepackage{amssymb}
\usepackage{amsthm}
\usepackage{enumitem}
\setlist[enumerate]{itemsep=0pt,topsep=0pt,parsep=0pt,leftmargin=0pt,itemindent=2.5em}
\setlist[itemize]{itemsep=0pt,topsep=0pt,parsep=0pt,leftmargin=0pt,itemindent=2.5em}
\usepackage{fancyhdr}
\pagestyle{fancy}
\fancyhf{}
\usepackage{graphicx}
\usepackage{hyperref}
\usepackage{indentfirst}
\usepackage{lmodern}


\begin{document}
\setlength{\abovedisplayskip}{1pt}
\setlength{\belowdisplayskip}{1pt}

\section{乘法原理}

乘法原理用于计算Descartes积的势. 为了区分Descartes积,
这里用$\displaystyle\bigotimes$表示对基数的求积.

\vspace{10pt}

\hypertarget{sec:定理1.1}{}
\noindent\textbf{定理1.1}

任给集合$A,B$, 有$|A\times B|=|A|\cdot|B|$.

\noindent{\kaishu 证明.}

记$A,B$的势分别为$\kappa,\lambda$, 则显然$A\times B\,\approx\,
\kappa\times\lambda\,\approx\,\kappa\lambda$. \hfill $\qedsymbol$

\vspace{10pt}

\hypertarget{sec:定理1.2}{}
\noindent\textbf{定理1.2}

任给有限集$X$和$X$上的映射$f$, 有$\displaystyle
\left|\prod_{x\in X}f(x)\right|=\bigotimes_{x\in X}|f(x)|$.

\noindent{\kaishu 证明.}

不妨设$X$非空, 取双射$\varphi:|X|\to X$. 对$n<|X|$归纳证明$\displaystyle
\left|\,\prod_{i=0}^nf\varphi(i)\,\right|=\bigotimes_{i=0}^n|f\varphi(i)|$.

首先, $n=0$时
\[
    \left|\,\prod_{i=0}^0f\varphi(i)\,\right|=\left|\,^1f\varphi(0)\right|=
    |f\varphi(0)|=\bigotimes_{i=0}^0|f\varphi(i)|\,;\\[3pt]
\]
其次, 假设命题对$n$成立且$n+1<|X|$, 那么易证
\vspace{3pt}
\[
    \prod_{i=0}^{n+1}f\varphi(i)\to\prod_{i=0}^nf\varphi(i)\times f\varphi(n+1),
    \quad p\mapsto(p\!\restriction_{n+1},p(n+1))\\[3pt]
\]
是双射, 从而
\vspace{3pt}
\[
    \left|\prod_{i=0}^{n+1}f\varphi(i)\right|=
    \left|\prod_{i=0}^nf\varphi(i)\times f\varphi(n+1)\right|=
    \bigotimes_{i=0}^n|f\varphi(i)|\cdot|f\varphi(n+1)|=
    \bigotimes_{i=0}^{n+1}|f\varphi(i)|,\\[6pt]
\]
归纳成立, 取$n=|X|-1$即得$\displaystyle\left|\prod_{x\in X}f(x)\right|=
\left|\prod_{i=0}^{|X|-1}f\varphi(i)\right|=
\bigotimes_{i=0}^{|X|-1}|f\varphi(i)|=\bigotimes_{x\in X}|f(x)|$.
\hfill $\qedsymbol$

\vspace{13pt} % 暂停.

\end{document}